\documentclass[10pt,a4paper]{amsart}
\usepackage[margin=19mm]{geometry}
\usepackage{amsmath,amssymb,mathtools}
\usepackage[scr=rsfs,cal=euler]{mathalfa}
\usepackage{xfrac}
\usepackage[unicode,hidelinks,bookmarks=false]{hyperref}
\newcommand{\ie}{i.e.\ }
\newcommand{\cf}{cf.\ }
\newcommand{\resp}{respectively\ }
\newcommand{\Subsection}{\subsection}
\providecommand{\nobreakdash}{\nobreak\hbox{-}}
\setlength{\emergencystretch}{2em}
\multlinegap0pt
\newtheorem{theorem}{Theorem}
\newtheorem{Conjecture}{Conjecture}
\newtheorem{prop}{Proposition}
\newtheorem{lemma}{Lemma}
\newtheorem{remark}{Remark}
\newcommand{\R}{{\mathbb R}}
\newcommand{\RV}{{\mathcal R}}
\newcommand{\T}{{\mathbf T}}
\newcommand{\I}{{\mathbf I}}
\newcommand{\RO}{{\RV\otimes\RV}}
\newcommand{\A}{{\mathbf A}}
\newcommand{\B}{{\mathcal B}}
\newcommand{\drm}{{\mathrm d}}
\newcommand{\M}{{\mathbf M}}
\begin{document}
\title[A short note on Meyers' theorem]{A short note on Meyers' theorem}
\author{M.A. Perelmuter}
\date{}
\maketitle
\noindent\emph{Source and licence.} A short note on Meyers' theorem. Version arXiv:2608.07030v1 (CC BY 4.0). This material is distributed under the Creative Commons Attribution 4.0 International licence, http://creativecommons.org/licenses/by/4.0/, as recorded in the arXiv OAI licence metadata for this exact version and shown on the abstract page https://arxiv.org/abs/2608.07030v1. Full rights notice: licenses/arxiv-2608.07030-CC-BY-4.0.txt.\par\smallskip
\noindent\textit{Editorial scope.} Counted unit: the original \texttt{prop} environment \texttt{eq:Cassese2} at source lines 101--110 together with its complete proof at lines 111--115; the delivered document keeps source lines 86--126 so that every referenced label and every citation resolves inside it. Native editable source is the paper's own concatenated TeX; the AMS wrapper is editorial formatting only. No publisher class, upstream script or unrelated artwork is redistributed.
\section{Main result}
Our main result is
\begin{theorem}\label{main} The Meyers interval contains
$$\left(\frac{26K-16}{13K-3}, \frac{26K-16}{13K-13}\right).$$
\end{theorem}
\begin{prop}\label{eq:Cassese} \cite[Proposition 4.1]{Cassese}
For $p\in(1,\infty)$ and $d\in {\mathbb N}\setminus\{0,1\}$,
\begin{equation}\label{RieszMatrix}
\|\RO\|_{L^p(\R^d,M_d(\R))\to L^p(\R^d, M_d(\R))}\le p^*-1,
\end{equation}
where $M_d(\R)$ is the space of $d\times d$ matrices over $\R$ and the matrix operator $\RO$ acts on a matrix-valued function $F\colon\R^d\to M_d(\R)$ by the rule
\begin{equation*}
((\RO)(F))_{i,l}=\sum\limits_j R_i R_j F_{j,l}.
\end{equation*}
\end{prop}

\begin{prop}\label{eq:Cassese2}
For $p\in(1,\infty)$ and $d\in {\mathbb N}\setminus\{0,1\}$,
\begin{equation}\label{RieszVector}
\|\RO\|_{p,p}\le p^*-1,
\end{equation}
where the operator $\RO$ acts on a vector-valued function $F\colon\R^d\to \R^d$ by the rule
\begin{equation*}
((\RO)(F))_{i}=\sum\limits_j R_i R_j F_{j}.
\end{equation*}
\end{prop}

\begin{proof}
To prove \eqref{RieszVector}, it suffices to consider the restriction of the operator in \eqref{RieszMatrix} to matrices with only one non-zero column and to use the fact that the subspace
$\Bigl\{F=(F_{ij}) \in L^p(\R^d,M_d(\R)):\;F_{ij}=0 \text{ for } j\neq 1\Bigr\}$
is an invariant subspace for the operator $\RO$.
\end{proof}

Now the triangle inequality gives
\begin{equation*}
\|\T\|_{p,p}\leq 2p^*-1.
\end{equation*}
The drawback of this estimate is that it does not capture the correct behavior for $p$ near $2$ ($\|\T\|_{2,2}=1$ due to the isometry of $\T$). But in the spirit of Lemma 9.1 from \cite{IS2001} we have the following
\begin{lemma}\label{InterpolLemma}
Let $1<p< \infty$, then
\begin{equation}\label{Interpol26}
\|\T\|_{p,p}\le 2.6(p^*-2)+1.
\end{equation}
\end{lemma}

\begin{thebibliography}{99}
\bibitem[Cassese]{Cassese} Cassese,~G.: \newblock Korn's inequality from the viewpoint of calculus of variations. \newblock arXiv:2603.22431.
\bibitem[IS2001]{IS2001} Iwaniec,~T., Sbordone,~C.: \newblock Quasiharmonic fields. \newblock Ann. I.H.Poincar\'e, Non Linear Analysis, \textbf{18}, No. 5, 519-572 (2001)
\end{thebibliography}
\end{document}
